\documentclass{article}
\usepackage{graphicx}
\usepackage[margin=1.5cm]{geometry}
\usepackage{amsmath}

\begin{document}
\twocolumn

\title{Warm Up: Energy I and Energy II}
\author{Prof. Jordan C. Hanson}

\maketitle

\section{Memory Bank}

\begin{itemize}
\item $v_f^2 = v_i^2 + 2 a \Delta x$ ... Kinematic equation without time
\item $W = \vec{F} \cdot \Delta \vec{x}$ ... Definition of work
\item $\vec{F} = k\Delta \vec{x}$ ... Hooke's Law, or the force of a spring
\item $W = \frac{1}{2}k(\Delta x)^2$ ... Work done to compress or stretch a spring by $\Delta x$
\item $KE = \frac{1}{2}m v^2$ ... Definition of Kinetic Energy
\item $W = KE_f - KE_i$ ... Work-energy theorem
\item $U = mgh$ ... Gravitational \textit{potential} energy, for a mass $m$ a height $h$ above the surface.
\item $R = v_i^2 \sin(2\theta)/g$ ... The range formula
\item $P = \Delta W/\Delta t$ ... Power is the ratio of work done to the time duration.  Units: J/s or Watts (W).
\end{itemize}

\section{Work, and the Work-Energy Theorem}

\begin{enumerate}
\item Suppose a spring is compressed by $\Delta \vec{x}$.  (a) If the spring ($k = 30$ N cm$^{-1}$) is compressed a distance of 5 cm to launch a 250 gram object, what is the kinetic energy as the object leaves the spring? (b) What is the speed of the object? (c) How high does the object rise, if it is shot straight upwards? \\ \vspace{4cm}
\item Now imagine the launcher from the previous exercies is tilted at a 45 degree angle.  What is the range? \\ \vspace{2cm}
\item Recall the kinematic equation without time (memory bank). (a) Rearrange the equation so that the velocity terms are on the left side, and divide both sides by 2.  (b) Multiply both sides by mass, $m$.  (c) Show that the right side is equal to the work done on $m$, and that we have proved a version of the work-energy theorem. \\ \vspace{4cm}
\end{enumerate}

\section{Power}

\begin{enumerate}
\item Suppose we push a 40 kg box along a smooth floor with $\mu_k = 0.05$ for a distance $40$ m.  (a) What is the work done? (b) If we perform this task in 30 seconds, what power did we generate? \\ \vspace{3cm}
\item The activity of cycling is 20 percent efficient.  That is, for every 100 Joules of energy we spend, we get 20 Joules of mechanical work.  (a) Suppose a cyclist is racing along at constant velocity.  The cyclist then climbs a 140 meter hill in 35 seconds.  The mass of the bicycle and cyclist is 70 kg.  What power was required to do this? (b) What total energy did the cyclist have to spend to do this?
\end{enumerate}

\end{document}
